% -------------------------------------------------------------------------
%  fig-coverage.tex  --  coverage monotonicity
%
%  Two panels.  Left: the covered region at time t, and a successor z sitting
%  inside it.  Right: after an update that RETIRES Z_1 -- the retired blob is
%  drawn dashed, strictly inside the blob that replaced it, so the eye sees
%  immediately that coverage grew rather than moved.  Labels are placed OUTSIDE
%  the blobs so nothing overlaps.
% -------------------------------------------------------------------------
\newcommand{\FigCoverage}{%
\begin{figure}[tb]
\centering
\begin{tikzpicture}[x=1mm,y=1mm]

% ---- panel: time t ------------------------------------------------------
\begin{scope}
  \draw[draw=black!35, line width=0.4pt, rounded corners=2pt]
    (0,-26) rectangle (56,10);
  \node[lbl, anchor=north west] at (1.5,9) {antichain $A_t$};

  \fill[black!10] (17,-11) ellipse (13 and 8);
  \draw[black, line width=0.5pt] (17,-11) ellipse (13 and 8);
  \node[lbl, anchor=north] at (17,-20.5) {$\aLU([Z_1])$};

  \fill[black!10] (42,-9) ellipse (10 and 6.5);
  \draw[black, line width=0.5pt] (42,-9) ellipse (10 and 6.5);
  \node[lbl, anchor=north] at (42,-17.5) {$\aLU([Z_2])$};

  \fill[black] (10,-13) circle (1.1);
  \node[lbl, anchor=east] at (8.2,-13) {$z$};
\end{scope}

% ---- transition ---------------------------------------------------------
\draw[fl] (59,-11) -- (69,-11);
\node[lbl, anchor=south] at (64,-10) {\code{add}};

% ---- panel: time t+1 ----------------------------------------------------
\begin{scope}[xshift=72mm]
  \draw[draw=black!35, line width=0.4pt, rounded corners=2pt]
    (0,-26) rectangle (56,10);
  \node[lbl, anchor=north west] at (1.5,9) {antichain $A_{t+1}$};

  \fill[black!10] (20,-10) ellipse (19 and 10);
  \draw[black, line width=0.5pt] (20,-10) ellipse (19 and 10);
  % the retired zone, drawn where it was: strictly inside its replacement
  \draw[black!55, line width=0.4pt, dash pattern=on 1.6pt off 1.4pt]
    (17,-11) ellipse (13 and 8);
  \node[lbl, anchor=north] at (20,-21.5) {$\aLU([Z_1 \sqcup Z_3])$};

  \fill[black!10] (42,-9) ellipse (10 and 6.5);
  \draw[black, line width=0.5pt] (42,-9) ellipse (10 and 6.5);

  \fill[black] (10,-13) circle (1.1);
  \node[lbl, anchor=east] at (8.2,-13) {$z$};
\end{scope}

\node[lblit, anchor=north, text width=124mm] at (64,-29)
  {$Z_1$ is retired --- but only into a zone that contains it (dashed).  The
   covered region grows; it never moves.  So $z$, once covered, stays covered.};

\end{tikzpicture}
\caption{Coverage monotonicity (\cref{thm:coverage}).  Concurrent workers
mutate an antichain while other workers read it, which is normally fatal to a
lock-free read.  It is safe here because of a structural property of the update
rule: a stored zone is retired only by subsumption or by absorption into a
convex hull, and in both cases the zone that replaces it \emph{contains} it.
The union of the covered regions is therefore non-decreasing over time, so a
``$z$ is already covered'' answer computed under a shared lock cannot be
invalidated by a writer that acquires the lock a nanosecond later.}
\label{fig:coverage}
\end{figure}}
